
% <<<PDFPAGE 37>>>
\chapter{长上下文与内存管理}

\section{注意力的平方复杂度}

注意力计算随序列长度 \(n\) 以 \(O(n^2)\) 增长，这是大语言模型长上下文推理的主要障碍。取 \mbox{\(n=1{,}000{,}000\)}，即 Gemini 1.5 Pro 等模型所用的百万词元上下文，朴素实现中的注意力矩阵需要 \(1{,}000{,}000^2\times 2~\text{字节}=2~\text{PB}\) 的内存，显然难以实现。即使 \(n=128{,}000\)（128K 上下文），注意力矩阵也占用 32~GB。

现代长上下文方法从几个方面处理这一问题。

\section{用于长上下文的 FlashAttention}

如第 7 章所述，FlashAttention 通过分块避免生成完整的注意力矩阵。计算复杂度仍为 \(O(n^2)\)，但 HBM 容量需求从 \(O(n^2)\) 降到 \(O(n)\)，使现有硬件可以处理含 128K 个词元的注意力计算。

\section{环形注意力}

环形注意力（Ring Attention，Liu 等，2023）把长上下文注意力计算分配到按环形连接的多台设备。每台设备持有序列的一部分 Q、K、V 张量，在计算局部注意力的同时，沿环传递 KV 块。借助在线 softmax，即使一台设备一次只接触一个 KV 块，仍可正确地归一化。设设备数为 \(p\)，环形注意力将每台设备的内存需求从 \(O(n^2)\) 降到 \(O(n/p)\)，从而能在设备集群上处理百万词元的上下文。

\section{上下文压缩}

上下文压缩（context compression）选择性地概括或删除信息量低的词元，从而省去对整段上下文计算完整注意力。

\textbf{LLMLingua}（Jiang 等，2023）用小型语言模型在给定上下文的条件下评估各词元的困惑度（perplexity），去掉困惑度低、容易预测的词元。许多任务中可实现 4–20 倍压缩，性能下降不到 5\%。

\textbf{AutoCompressor} 微调模型本身，使其将长上下文概括为固定数量的摘要向量（summary vectors），再将这些向量用作压缩后的 KV 缓存。

\section{作为线性注意力的状态空间模型}

Mamba（Gu 与 Dao，2023）及其后续模型属于状态空间模型。它们以学习到的循环状态代替注意力来处理序列，计算方式为
\[
  h_t=Ah_{t-1}+Bx_t,\qquad y_t=Ch_t.
\]
% <<<PDFPAGE 38>>>
其中，选择性状态空间模型（selective SSM，S6）的 \(A\)、\(B\) 和 \(C\) 取决于输入。这使推理耗时为 \(O(n)\)，内存占用为 \(O(1)\)，因为隐状态 \(h\) 的大小不随序列长度变化。在推理时，SSM 等效于循环神经网络（Recurrent Neural Network，RNN）。每收到一个新词元，只需更新状态即可，无须保存 KV 缓存。代价是，当任务需要准确回忆前文某处内容时，SSM 从上下文中检索信息的能力弱于注意力机制。

\begin{interview}
\question{1}{环形注意力如何实现分布式长上下文推理？通信开销有多大？}
环形注意力将序列分配给 \(p\) 台设备。设备 \(i\) 持有查询块 \(Q_i\) 和键值块 \((K_i,V_i)\)。在 \(p\) 轮通信中，每一轮设备都将手头的 KV 块发送给环中的下一台设备，同时接收上一台设备发来的 KV 块。等待下一块数据的时间可与计算重叠。各设备计算本地 \(Q_i\) 与当前 KV 块之间的部分注意力分数，并累积在线 softmax 的统计量。经过 \(p\) 轮，每台设备都得到了其查询块对应的正确注意力输出。

\textbf{通信开销。}每轮传输 \(n/p\) 个词元的 KV 缓存，即 K、V 两个张量。按每台设备计，每层传输总量为 \(n\times d_{\mathrm{head}}\times n_{\mathrm{heads}}\times 2\times\text{精度}\)，恰好等于整层的 KV 数据量，没有额外的通信开销。通信与计算采用流水线方式并行执行。

\question{2}{请比较 Mamba（SSM）和 Transformer 注意力在长上下文推理中的表现。你会分别在什么情况下选用它们？}
\textbf{Transformer 注意力}需要 \(O(n^2)\) 的计算量和 \(O(n)\) 的 KV 缓存空间。它擅长从上下文的任意位置准确回忆某项细节，例如“这份文档第 3 页怎样描述 X？”采用 FlashAttention 后，性能随上下文增长而逐渐下降。

\textbf{Mamba} 的计算量为 \(O(n)\)，状态空间为 \(O(1)\)，因为其循环状态大小固定。它处理极长序列时效率很高，每个词元的计算量恒定，不随上下文总长度变化。它的弱点是难以从相隔很远的上下文中准确检索信息，在依赖检索的评测（如 RULER 和口令检索，passkey retrieval）中表现不及注意力。

\textbf{适合 Mamba 的场景}包括超长文档的流式推理、时间序列预测，以及需要跟踪状态演变、较少检索具体事实的应用。Jamba、Zamba 以及 Mamba-2 与注意力结合等\textbf{混合模型}兼用两种方法，让 SSM 层承担大部分处理，仅配置少量注意力层用于检索。

\question{3}{LLMLingua 如何压缩提示？提示压缩存在哪些风险？}
LLMLingua 用小型代理大语言模型，在给定某个词元的上下文（其中可能已经做过压缩）的条件下估计其困惑度。困惑度低的词元容易预测，因而被视为信息量较低的删除候选。算法分三步。（1）在粗粒度上确定保留哪些句子或短语。（2）在细粒度上删除保留片段中困惑度低的词元。（3）根据目标压缩比决定删减程度。

\textbf{风险。}（1）代理模型与目标模型对“可预测”的判断可能不同，导致重要词元被删掉。（2）删除词元会改变后续所有词元所在的位置，采用绝对位置编码的模型尤其可能受影响。（3）对依赖检索的任务，例如准确查找姓名或数字，即使困惑度低的词元也可能至关重要。
% <<<PDFPAGE 39>>>
（4）压缩不可逆，删错后无法恢复。实际部署前，应在目标应用的任务分布上检验 LLMLingua 的效果。

\question{4}{YaRN 和 NTK 感知缩放在扩展 RoPE 上下文时有什么区别？}
这两种方法都试图将采用 RoPE、在上下文长度 \(L\) 上训练的模型扩展到更长的上下文 \(L'>L\)，同时避免微调或只做少量微调。\textbf{NTK 感知缩放}（bloc97，2023）重新缩放 RoPE 的频率基数，将 \(b=10{,}000\) 换成 \(b'=b\cdot(L'/L)^{d/(d-2)}\)。这让旋转角在扩展范围内分布得更均匀，避免高频维度的旋转角绕回。它对所有频率维度采用同一种缩放方式。

\textbf{YaRN}（Peng 等，2023）的处理更细致。它保持低频维度不变，因为这些维度已经覆盖了较远的范围，对中频维度采用 NTK 式缩放，对高频维度做线性插值。这样既可保留短序列上与训练分布一致的行为，又能处理更长的序列。YaRN 还进行注意力温度校正（attention temperature correction），将注意力分数乘以 \(\sqrt{1/t}\)，以补偿注意力分布的变化。在扩展上下文长度后的困惑度和检索评测中，YaRN 的表现持续优于 NTK 感知缩放。

\question{5}{RAG 是什么？它与长上下文推理有何关系？你何时会选用哪一种？}
检索增强生成（Retrieval-Augmented Generation，RAG）在接到查询时，从向量数据库中选取相关文档，只把最相关的段落插入提示。它借助检索系统免去了使用极长上下文窗口的需求。相较之下，\textbf{长上下文推理}把所有可能相关的信息直接放入模型的上下文窗口。

\textbf{适合 RAG 的情形}包括（1）语料库极大，例如包含数百万份文档，远超可用的上下文窗口。（2）检索足够准确，查询明确指向特定文档。（3）索引检索的时延可以接受。（4）计算长上下文的成本过高。

\textbf{适合长上下文推理的情形}包括（1）任务需要综合或比较文档多个部分的内容，例如对整本书进行复杂推理。（2）检索结果噪声较多，完整上下文有助于模型分辨信息是否相关。（3）语料库足够小，例如不足 100 万词元，可以整个装入上下文。理想的系统往往组合两者，先用 RAG 缩小候选范围，再用长上下文推理分析检索出的段落。
\end{interview}
